% Chapter 15 — Testability and the Executable-Check Regime (WILSON). Budget 5.5 pp.
% Discharges Constitution v1.2 §13 (two layers of correctness) and the §2 Testability commitment.
\section{Testability and the Executable-Check Regime}\label{ch:testability}

This chapter discharges constitution clauses C-13.1--C-13.3 (the two layers of correctness) and the Testability commitment C-2.5 (jointly with Chapter~\ref{ch:objective}). The claim it makes good is the constitution's: the
ledger is correct because it is built to be tested. Every invariant catalogued in
Chapter~\ref{ch:invariants} is an executable check run over generated products, events, and
histories, never a property defended only in prose; and correctness composes, because the
steps are pure and their composition is fixed --- the map, then the fold --- so verifying
each part verifies the whole. This chapter does not restate the invariants. It makes them
the surface a test regime exercises, states the regime, and turns the specification's three
binding conditions into executable gates rather than intentions.

\subsection{Two layers of correctness}

Correctness rests on two layers, and each is checked by its own means.

\emph{Structural} correctness holds by construction, at the door. An illegal state is not
detected and rejected; it is unrepresentable --- no program writes it. A balance changes
only through a move (Chapter~\ref{ch:objects}); a marker coordinate is unwritable without
its governing agreement reference (Chapter~\ref{ch:collateral}); a knocked note has no
un-knock constructor (Chapter~\ref{ch:objects}, Invariant~\ref{inv:graph-consistency});
conservation is automatic because every move is paired-leg and the system is closed. These
properties hold whatever the contract that proposed the transaction did --- correct,
careless, or malicious --- because the Transaction Executor admits on structure alone and
is the sole writer (Chapter~\ref{ch:machines}). The test obligation for this layer is
negative: no generated event, however adversarial, drives the ledger into a state the type
discipline forbids or the door should refuse.

\emph{Economic} correctness is not guaranteed at admission; it is made checkable by
recomputation. A contract may emit a perfectly well-formed transaction that books the wrong
economics. The check is to re-run the map on the recorded event and compare the result with
what was recorded.

\begin{lstlisting}
-- the economic oracle: re-run the map on the recorded event and compare
prop_recompute :: Recorded (Event, Transaction) -> Property
prop_recompute (e, txRecorded) = contract e (ledgerBefore e) === txRecorded
-- a mismatch is a defect of the CONTRACT, not the ledger; it is repaired FORWARD by an
-- authorised compensating transaction through the same door -- never by an edit to history.
\end{lstlisting}

Purity is what makes this oracle exist at all: a contract reads only the event and the
record (Chapter~\ref{ch:contracts}), so its recomputation is exact and unconditional. A
mismatch is a defect of the contract, never of the ledger, and it is repaired forward the way every
money change is made --- by an authorised compensating transaction through the same door
(the forward-repair discipline below). The regime-bit
repair of Chapter~\ref{ch:collateral} (a new terms version \emph{plus} the authorised compensating
rebooking of the affected interval) is the worked instance; the valuations of
Chapter~\ref{ch:valuation} are recomputable under the same discipline. The two layers
divide the labour exactly and neither substitutes for the other: structure is proved once
and tested negatively, economics is re-derived per recorded event and tested by comparison,
and a structurally valid transaction can still be economically wrong.

\paragraph{Forward repair: the correction event and the authorisation gate.}
A wrong recorded event --- a wrong price, a wrong payoff --- is repaired forward, never
edited (constitution C-12.4). The discovery is itself a recorded \emph{correction event}
that names the wrong event's identifier, admitted through the one door like any event. It
mints a visible, deadline-bearing \emph{open item} --- a deadline, a discharge predicate,
and a fallback, the obligation pattern of Chapter~\ref{ch:invariants} --- so the discrepancy
stands on the record until repaired. Two things then part company. A \emph{view} recomputes
automatically: the as-known-at projection stands, the as-of projection recomputes over the
corrected prefix (Theorem~\ref{thm:timetravel}). \emph{Money} does not: because money already
moved cannot be unsent, a compensating transaction that repairs a balance is agreed with the
counterparty and authorised by a person before it fires. The authorisation is a recorded
observation; the compensating transaction references both the correction event and its
authorisation, or the door refuses it. In phase~0 no correction fires without human
intervention.

\begin{lstlisting}
data Correction    = Correction    { wrongEventId :: EventId,  reason   :: Text }
-- the discovery: a moveless recorded event, through the one door, naming the wrong event.
data Authorisation = Authorisation
  { corrects :: EventId, approver :: PersonId, counterpartyAssent :: ObservationId }
-- a recorded approval carrying BOTH C-12.4 conditions: agreed with the counterparty
-- (the assent, a recorded observation) and authorised by a person (the approver).
compensating :: CorrectionId -> AuthorisationId -> [Move] -> Transaction
-- door-checkable definitions: isCompensating = carries a Correction (wrongEventId)
-- reference; hasAuthorisationRef = references an Authorisation whose corrects field
-- names that wrongEventId, carrying approver AND counterpartyAssent.
-- the guarantee is the DOOR's, not the constructor's -- an untrusted proposer can
-- bypass `compensating` -- so the property asserts the refusal:
prop_noRepairWithoutAuth :: Transaction -> Property
prop_noRepairWithoutAuth tx =
  isCompensating tx && not (hasAuthorisationRef tx) ==> refusedAtDoor tx
genUnauthorisedRepair :: Gen Transaction  -- ADVERSARIAL: reference a correction, omit
  -- the authorisation, bypass the constructor -- as the untrusted Executor could.
prop_noRepairAuth_fires =                 -- the refusal witness, asserted, not hoped
  checkCoverage $ cover 1.0 unauthorisedRepair
    "correction-referencing repair without authorisation -> refused" genUnauthorisedRepair
\end{lstlisting}

A view is free because it stores nothing and settles nothing; a compensating transaction
moves owned mass, so it passes the same human gate every real move passes.

\subsection{The generator universe and the adversarial scheduler}

The enumerations this specification builds on are genuinely closed. On the ledger's side:
the product-graph node and edge sets --- the lifecycle states and guarded transitions ---
fixed at registration (Chapter~\ref{ch:objects}); the collateral regimes
\{title-transfer, pledge, STM/CTM\} and the coordinate planes \{owned, lent, posted\}
(Chapter~\ref{ch:collateral}); the component dimensions of the data-kind registry and the
boundary failure modes W1--W4 (Chapter~\ref{ch:marketdata}); and the trigger and event
kinds of a contract (Chapter~\ref{ch:contracts}). On the CDM side the genuinely closed
enumerations are the day-count fraction, the option type (call or put), and the
primitive-instruction set (Chapter~\ref{ch:cdm}). These are closed choices, and they seed
the generator universe; a co-populated, cardinality-many composite is not a closed choice
and is not one of them. Because the closed enumerations are closed and declared, one
generator universe ranges over the whole domain, and its coverage of each enumeration is
itself a checkable fact --- every inhabitant is reachable by the generators, so a property
that never exercises a case is a visible gap, not a silent one.

\begin{lstlisting}
genProduct :: Gen ProductGraph          -- a well-typed graph over the closed node/edge set
genEvent   :: ProductGraph -> Gen Event -- on-menu guards, and off-menu adversarial events
genHistory :: Gen [Transaction]         -- interleaved admitted walks of many products
-- every generator is paired with a shrinker: a violation shrinks to a minimal history,
-- and every counterexample is replayable from its seed (a defect is a reproduction).
\end{lstlisting}

This generator regime is the simulated ledger of Chapter~\ref{ch:virtual}
(Section~\ref{sec:simulated}) run under the test measure: the same branch-point,
seed, and one-door discipline, aimed at refutation rather than pricing.

Generation is adversarial by design, because the test environment must be harder than
production. The Event Monitor and the Events Executor are outside the trust boundary and may
duplicate, delay, and reorder firings; integrity owes them nothing
(Chapter~\ref{ch:machines}). The scheduler therefore explores every legal interleaving the
Transaction Executor's total order permits --- duplicated proposals, late firings, events
delivered out of arrival order --- and injects pathological-but-legal behaviour rather than
only the happy path: this is precisely what exercises idempotence and the temporal model
below. One discipline guards the generators against vacuity: a property whose precondition
is never generated passes without witnessing anything. A generated history that never knocks
a pledged note, never re-posts received collateral, never delivers a duplicate is not
evidence; the generators are measured by the states they reach, not the cases they cover.
Throughout this chapter each property's firing witness asserts its branch at $\ge 1\%$ of
generated histories under \texttt{checkCoverage}; an unwitnessed branch is a \emph{failed}
run, not a silent green, by \S 2's zero-firings rule. This standing gate is not restated per
property.

\subsection{The test surface: the invariant catalogue and the product graph}

The invariants of Chapter~\ref{ch:invariants} are the test surface, in a ramp from the
universal to the domain-specific. At the base sit the properties no history may violate: no
transaction half-applies (atomicity), no balance is created or destroyed off the paired-leg
move (conservation per (unit, coordinate, agreement)), the log is append-only and
hash-chained, and every arrow is deterministic. Above them sit the structural invariants --- consistency of
reference, writer discipline --- and above those the safety properties of the collateral
planes: coverage is an invariant, checked at the door on every admission, while sufficiency
is an obligation, lawfully false intraday and bounded by a deadline, and the test regime
must assert the first as an invariant and the second only as a liveness obligation. Stating
sufficiency as an invariant would be a false theorem the generators refute at once, and
insisting on it would forbid lawful states: the two guards are never collapsed, in the tests
least of all.

\paragraph{Conservation at agreement grain, and the netting collapse it catches.}
Conservation is checked at the finest grain the coordinate vector has --- per (unit,
coordinate, agreement), Theorem~\ref{thm:conservation}. The teeth are the cross-agreement
history: post 100 under $G_1$, receive 100 under $G_2$, aggregate posted axis zero. A
property reading only the aggregate sees zero and passes \emph{vacuously}, reporting a
fully encumbered wallet as unencumbered (Chapter~\ref{ch:reporting}); the finer law keeps
$+100$ and $-100$ as two coordinates, and the gross-encumbrance property reads them as
$100$, not $0$.

\begin{lstlisting}
prop_conservation :: History -> Unit -> Coord -> Agreement -> Property  -- Theorem 14.1
prop_conservation h u c g = sum [ balance c g w u (fold h) | w <- wallets h ] === 0

prop_encumbranceGross :: History -> Wallet -> Unit -> Property          -- Ch. 12, gross
prop_encumbranceGross h w u =
  reportedEncumbrance h w u === sum [ max (postedG h w u g) 0 | g <- agreements h ]
  -- the signed net across agreements is NEVER the encumbrance figure.

genCrossAgreement :: Gen History  -- draw a netting-flag with rate r; when it fires,
  -- post q of u under G1 (bal_posted,G1 = +q) and receive q of u under G2 (= -q),
  -- so the AGGREGATE posted axis is 0 while the two coordinates are +q and -q.
prop_netting_fires =                     -- the firing witness, asserted, not hoped
  checkCoverage $ cover 1.0 isCrossAgreementNetting
    "post under G1 and receive under G2 (aggregate posted = 0)" genCrossAgreement
\end{lstlisting}

The netting branch is a genuine precondition, not a tautology: a generator that never posts
and receives the same unit under two agreements would satisfy \texttt{prop\_encumbranceGross}
without ever witnessing the collapse it exists to catch. \texttt{genCrossAgreement} constructs
the branch by construction, and its firing witness asserts it (\S 2). On every history that
clears the precondition, \texttt{prop\_conservation} fires on each of $G_1$ and $G_2$ and
finds each coordinate conserved, and \texttt{prop\_encumbranceGross} fires and finds the
reported encumbrance $100$ where a net-reading report would show $0$ --- the single-axis
property passing vacuously while the agreement-grain pair catches the collapse.

\paragraph{The coverage sign convention, and the over-re-post it refuses.}
The rehypothecation bound nets the signed posted coordinates against possession
(Invariant~\ref{inv:coverage}), and the whole bound rests on one sign convention: on
the posted coordinate collateral posted \emph{out} is positive and collateral
received \emph{in} is negative --- received-negative. Get the sign backwards and the
bound inverts: a wallet that received collateral would appear able to re-post it
without limit. The property asserts the convention by asserting the refusal it
forces. Take a wallet with owned $0$; it receives $100$ under one agreement (posted
coordinate $-100$) and attempts to re-post $200$ under another ($+200$). The net is
$200-100=+100>\max(0,0)=0$, so the admission is refused at the door; under the
inverted sign the same net would read $-100\le 0$ and admit --- the exact defect the
property catches. The precondition --- a zero-owned wallet that receives and then
re-posts \emph{more} than it received --- is a genuine one, not a tautology.

\begin{lstlisting}
prop_coverageSignConvention :: History -> Wallet -> Unit -> Property  -- Invariant 14.6
prop_coverageSignConvention h w u =
  admitted h ==> coverageNet h w u <= max (ownedBal h w u) 0
  where coverageNet h w u = sum [ postedG h w u g | g <- agreements h ]
  -- received IN contributes NEGATIVE, posted OUT contributes POSITIVE.

genOverRepost :: Gen History  -- draw a re-post-flag with rate r; when it fires, take a
  -- wallet with owned 0, receive 100 of u under G1 (bal_posted,G1 = -100), then attempt
  -- to re-post 200 of u under G2 (bal_posted,G2 = +200): net +100 > max(0,0)=0, REFUSED.
prop_coverageSign_fires =                 -- the refusal witness, asserted, not hoped
  checkCoverage $ cover 1.0 isOverRepostOnZeroOwned
    "owned 0, received 100, re-post 200 -> net +100, refused" genOverRepost
\end{lstlisting}

When the re-post-flag fires the generator constructs the over-re-post by construction
--- owned $0$, received $100$, re-post attempt $200$ --- and its refusal witness asserts
the branch (\S 2). Flip the convention to received-positive and the refused net $+100$ reads
$-100 \le 0$ and admits: \texttt{prop\_coverageSignConvention} reports exactly that
inversion, while the lawful chain of Invariant~\ref{inv:coverage} --- received $100$,
re-post $60$, net $-40 \le 0$ --- stays admitted.

Carried onto this surface are the four properties of the product-graph consistency invariant
(Invariant~\ref{inv:graph-consistency}, defined in Chapter~\ref{ch:objects}), stated here as
executable checks over generated products and histories:

\begin{lstlisting}
-- for every unit, at every commit, the three interpreters of its graph agree:
prop_watchesAreOutEdges   :: History -> Property  -- (i) watch list == out-edges(current node)
prop_firedMatchesOneEdge  :: History -> Property  -- (ii) fired event -> exactly one edge, lands on its target
prop_actionEqualsRecorded :: History -> Property  -- (iii) emitted tx == edge action applied to recorded inputs
prop_traversalExpiresSibs :: History -> Property  -- (iv) traversal expires exactly the declared siblings
\end{lstlisting}

Property~(i) is checked at every log position: the Event Monitor's live watch set is folded
from the record and compared with the out-edge set of the unit's current node. Property~(ii)
rejects any firing that matches no edge or more than one, or that lands the unit off the
edge's declared target. Property~(iii) is the graph-level face of the economic oracle above:
the transaction the contract emitted must equal the edge's declared action applied to the
recorded inputs. Property~(iv) walks the sibling set and confirms that a traversal expires
those edges and no others. The graph declares watches, state schema, and actions from one
object; these four properties are the statement, made executable, that its three interpreters
never disagree.

\paragraph{Term adjustment on a corporate action, and the phantom knock it refuses.}
Across a split of a referenced underlying, the corporate-action edge
(Section~\ref{sec:termsheet-graphs}) appends a new ProductTerms version through the
term-adjustment operator (Chapter~\ref{ch:marketdata}), and it carries its own executable
witness (Invariant~\ref{inv:consistency}): the payoff read on the appended terms equals the
payoff on the pre-event terms, and no guard fires that the pre-event terms would not have
fired. The precondition --- a generated derivative \emph{and} a generated corporate action
on its referenced underlying --- is genuine, not a tautology.

\begin{lstlisting}
prop_termAdjustmentInvariance :: History -> Property
prop_termAdjustmentInvariance h = hasCAonUnderlying h ==>
       payoffOn (postEventTerms h) (postEventFrame h)
   === payoffOn (preEventTerms  h) (preEventFrame  h)      -- payoff invariant across the CA
  .&&. not (spuriousGuardFires h)  -- no edge fires that the pre-event terms would not

genDerivWithCA :: Gen History
-- 1. draw a derivative whose terms REFERENCE an underlying (OT-1's barrier on OMEGA);
-- 2. draw a CA-flag with rate r; when it fires, emit a two-for-one split of that referenced
--    underlying inside the derivative's live window, then record the next underlying close
--    in the GAP between the adjusted and the stale barrier (e.g. 47.50, with 80.00 -> 40.00)
--    -- the exact phantom-knock witness the edge must refuse.
prop_termAdj_fires =                      -- the firing witness, asserted, not hoped
  checkCoverage $ cover 1.0 hasCAonUnderlying
    "split on a referenced underlying, next close in the phantom gap" genDerivWithCA
\end{lstlisting}

When the CA-flag fires the generator constructs the split and prints the next close in the
phantom gap, and its firing witness asserts the branch (\S 2). On every history that clears
the precondition the payoff is invariant across the appended terms and the stale-barrier
knock --- $47.50 \le 80.00$, a payout on a stock that never fell --- does not fire, because
the live edge reads the adjusted $40.00$ and $47.50 > 40.00$. Strip the corporate-action edge
and the barrier stays $80.00$, the knock fires, and \texttt{prop\_termAdjustmentInvariance}
reports the phantom the edge exists to refuse.

\paragraph{The settlement-obligation graph, and the fail branch it must walk.} The
settlement-obligation unit (Chapter~\ref{ch:settlement}) is a product graph like any
other, tested by the four properties above (Invariant~\ref{inv:graph-consistency}) --- no
new machinery. What the generator must not leave unwitnessed is the off-happy-path walk:
the fail edge, the cum record-date branch, and the fourth-quadrant mirror branch, each
constructed below; a generator that only ever settles would satisfy the four properties
without exercising any of them.

\begin{lstlisting}
genSettlementHistory :: Gen History
-- 1. instruct a trade -> mint the settlement-obligation unit at node `instructed';
-- 2. draw a fail-flag with rate r; when it fires, deliver a settlement-fail event so the unit
--    walks instructed -> failed (the SAME unit at its fails-cascade node, never a fresh one);
-- 3. draw a cum record-date flag with rate r; drop a record date inside the instruction-to-
--    settlement gap on a CUM trade, so the unit raises the market-claim leg (deliverer -,
--    cum-buyer +) at its `instructed' node;
-- 4. draw an ex-settled flag with rate r; settle an EX trade BEFORE the record date, so the
--    register shows the buyer and the unit, at its `settled' node, raises the MIRROR
--    market-claim leg (ex-buyer -, seller +) -- the fourth quadrant, reachable under a
--    due-bills regime whose ex-date falls after the payment date.
prop_settlementFail_fires =                 -- the firing witness, asserted, not hoped
  checkCoverage $ cover 1.0 reachesFailedNode
    "settlement-obligation unit walks instructed -> failed" genSettlementHistory
prop_mirrorClaim_fires =                     -- the fourth-quadrant firing witness, asserted
  checkCoverage $ cover 1.0 raisesMirrorClaim
    "ex trade settled before record date raises the mirror market-claim (buyer -> seller)"
    genSettlementHistory
\end{lstlisting}

When the fail-flag fires the generator walks the unit to its \texttt{failed} node by
construction, so \texttt{prop\_firedMatchesOneEdge} confirms the fail event lands it there
and \texttt{prop\_actionEqualsRecorded} confirms the emitted transaction is the edge's
declared action --- the fails-cascade obligation, never a reversal of history. The firing
witness asserts the fail branch (\S 2).

The fourth quadrant needs its own witness. A generator that only ever leaves an ex trade
unsettled at the record date --- the coherent-calendar case --- raises no mirror leg and
never exercises the ex-settled corner. Branch~4 settles an ex trade before the record date
by construction, so the record-date firing reads the buyer as registered holder against an
ex entitlement and raises the mirror market-claim leg from buyer to seller
(Chapter~\ref{ch:settlement}); \texttt{prop\_mirrorClaim\_fires} asserts
\texttt{raisesMirrorClaim} (\S 2). The branch
models the lawful due-bills regime, where an ex trade can settle before the record date, and
the property confirms the distribution routes to the entitled seller on the record, not to
the registered buyer --- the reflection, across the settlement node, of the cum-unsettled
market claim.

\paragraph{The position-state collection: two entitlements in flight, each retiring alone.}
PositionState holds a collection keyed by the cause-derived identifier
(Chapter~\ref{ch:homes}): two entitlements can be live on one (wallet, unit) pair at once,
each retiring exactly when its own obligation discharges. A generator that declares only one
dividend at a time would satisfy every balance property without ever placing two facts in
one slot.

\begin{lstlisting}
prop_twoEntitlementsRetireIndependently :: History -> Property
prop_twoEntitlementsRetireIndependently h = twoDividendsInFlight h ==>
      -- after both record dates, the pair carries BOTH entitlements, keyed distinctly:
       Map.size (posFacts h alpha acme) === 2
      -- paying the first retires exactly its entry; the second stays live to its own date:
  .&&. posFacts (afterPay1 h) alpha acme === Map.singleton (txid2 h) (Owed d2)
  .&&. posFacts (afterPay2 h) alpha acme === Map.empty

genTwoDividends :: Gen History  -- draw a two-dividend flag with rate r; when it fires, declare
  -- a dividend (record r1, pay p1) and a SECOND (record r2, pay p2) with r1<r2<p1<p2, so both
  -- are owed at once and their payment dates are distinct -- the two-in-flight branch.
prop_twoInFlight_fires =                    -- the firing witness, asserted, not hoped
  checkCoverage $ cover 1.0 twoDividendsInFlight
    "two dividends owed to one holder at once, distinct payment dates" genTwoDividends
\end{lstlisting}

When the flag fires the generator overlaps the two dividends by construction ($r_1<r_2<p_1<p_2$),
so between $p_1$ and $p_2$ exactly one entitlement is live, and its firing witness asserts the
branch (\S 2). On every history that clears the precondition, both
entitlements are keyed distinctly and live together, paying the first retires its entry alone,
and the second discharges on its own payment date. The property fails on any design that keeps
one slot per pair: the second record-date firing would overwrite the first, and the collection
could never hold the two the routine Tuesday puts in flight.

\paragraph{The coordinated cascade: delivered twice, each leg once, all $n$ committed.}
One cause can move many units: a cascade of $n$ legs sharing one $\mathit{causeEventId}$,
each with its own $\mathit{txid}$ (the derivation rule of Chapter~\ref{ch:machines}).
The property asserts both directions --- no leg dropped, no leg duplicated --- by
delivering the whole cascade twice, reordered.

\begin{lstlisting}
prop_cascadeIdempotence :: History -> Property
prop_cascadeIdempotence h = isCascade h ==>
      -- ALL n commit: the committed identifiers are exactly the cascade's n legs
      Set.fromList (committedTxids h) === Set.fromList (cascadeTxids h)
      -- EACH exactly once: the second delivery adds nothing and drops nothing
  .&&. length (committedTxids h) === n
  where n = length (cascadeTxids h)

genCascade :: Gen History  -- draw a cascade-flag with rate r; when it fires:
  -- 1. one cause on m>=2 units (a corporate action adjusting several referencing units)
  --    -> m legs sharing ONE causeEventId, distinct unitId (seq splits any same-unit legs);
  -- 2. deliver the m proposals once           -> n=m distinct txids, all commit;
  -- 3. deliver the SAME m proposals AGAIN, REORDERED (duplicate + reorder)
  --    -> each recomputes its identical txid; the door commits it zero further times.
prop_cascadeIdempotence_fires =            -- the firing witness, asserted, not hoped
  checkCoverage $ cover 1.0 isCascade
    "n>=2-leg cascade delivered twice, reordered" genCascade
\end{lstlisting}

When the cascade-flag fires the generator builds an $m\ge 2$-leg cascade and delivers
it a second time in permuted order, and its firing witness asserts the branch (\S 2). On
every history that clears the
precondition, the committed identifiers are exactly the cascade's $n$ distinct
$\mathit{txid}$s --- \emph{all $n$ commit}, because the legs differ in
$(\mathit{unitId},\mathit{seq})$, and \emph{each exactly once}, because the reordered
second delivery recomputes identifiers already on the log. Key idempotence on
$\mathit{causeEventId}$ alone and legs drop; drop the retry check and every leg
double-commits --- the property fails on both broken designs.

\paragraph{Independent transactions commute; dependent ones do not.}
Two adjacent admitted transactions with disjoint footprints swap to bit-identical state;
two that share a footprint element may not (Proposition~\ref{prop:independent-commute}).
The pair is tested both ways, and the second half keeps the first honest: without a
witnessed non-commuting case, ``swapping preserves state'' could pass on a generator that
only ever draws independent pairs. The first property swaps a disjoint pair and asserts the
full state is identical; the second swaps a buy against a record-date event on the
\emph{same} stock and wallet and asserts the two orders differ, because the entitlement
differs.

\begin{lstlisting}
prop_independentSwapInvariant :: History -> Property  -- Prop. 14 (independent commutation)
prop_independentSwapInvariant h = hasDisjointAdjacentPair h ==>
      -- swap the two adjacent txs; full state (balances + three homes) is bit-identical.
      -- oracle is fullState . fold, independent of the swap under test (NOT x===x):
      -- one side folds the given order, the other folds the swapped order.
  fullState (fold h) === fullState (fold (swapAdjacent h))

genDisjointPair :: Gen History  -- draw a disjoint-pair flag with rate r; when it fires,
  -- build two adjacent txs whose footprints share NOTHING, by construction: draw two
  -- DIFFERENT units and two DIFFERENT wallets, so the (w,u) coverage read-fibres are
  -- wholly disjoint (Prop. 14: the footprint includes the door's reads at fibre grain
  -- -- different agreements on one (w,u) would NOT suffice) and the three homes' targets
  -- differ.
prop_independentSwapInvariant_fires =    -- the firing witness, asserted, not hoped
  checkCoverage $ cover 1.0 hasDisjointAdjacentPair
    "adjacent pair, disjoint across all four change kinds" genDisjointPair

prop_dependentSwapDiffers :: History -> Property             -- the C-2.7 witness
prop_dependentSwapDiffers h = hasBuyThenRecordDate h ==>
      -- buy ACME then meet record date -> entitled; meet record date then buy -> not.
      -- the two orders share the (wallet, ACME) footprint, so they must NOT commute:
  entitlement (fold h) =/= entitlement (fold (swapAdjacent h))

genBuyRecordDate :: Gen History  -- draw a dependent-pair flag with rate r; when it fires,
  -- place a buy of ACME into W-ALPHA immediately adjacent to ACME's record-date event,
  -- SAME stock, SAME wallet -> shared footprint -> the entitlement depends on the order.
prop_dependentSwapDiffers_fires =        -- the non-commuting firing witness, asserted
  checkCoverage $ cover 1.0 hasBuyThenRecordDate
    "buy and record-date on the same stock and wallet, adjacent" genBuyRecordDate
\end{lstlisting}

Both preconditions are genuine, not tautologies, and the two guard against opposite
failures. \texttt{genDisjointPair} constructs disjointness by drawing two different units,
wallets, and agreements, so every generated pair shares no footprint element and the swap
is provably harmless; a generator that never produced such a pair would let
\texttt{prop\_independentSwapInvariant} pass without witnessing a single lawful swap.
\texttt{genBuyRecordDate} constructs the
opposite --- a buy and a record-date event on one stock in one wallet --- so the shared
footprint makes the entitlement order-dependent; its oracle, \texttt{entitlement . fold},
is computed independently on each order rather than compared to itself, so a green run means
the two orders genuinely diverged. Each firing witness asserts its branch (\S 2). Together the
pair pins C-2.7 executably: the harmless swap is bit-identical, the dependent swap is not,
and neither branch is vacuous.

\subsection{The three binding conditions as executable gates}

Three conditions the specification relies on are stated here as gates a candidate
implementation must pass, not as intentions it may assert.

\paragraph{(i) Moves are the sole balance mutator.}
The constitution makes the move the only operation that changes a balance
(Chapter~\ref{ch:objects}); the gate makes that mechanically decidable. Applying a
transaction, then applying only its move list, must leave identical balances --- so a
unit-state, position-state, or terms change that touched a balance is caught as a divergence.

\begin{lstlisting}
prop_movesOnly :: Transaction -> Property
prop_movesOnly tx = balances (apply tx l0) === balances (apply (onlyMoves tx) l0)
  where l0 = ledgerBefore tx
-- corollary, tested jointly:  apply (dropMoves tx) l0  leaves every balance fixed.
\end{lstlisting}

\paragraph{(ii) Determinism and replay.}
The clock is the single input not on the record, and the Event Monitor is the single place it
is consulted (Chapter~\ref{ch:machines}); every arrow downstream is a function of the record
alone. Two gates follow. \emph{Idempotence:} a duplicated proposal, carrying the same
cause-derived identifier, commits once and produces the identical transaction --- a
duplicate is inert. A late or out-of-order proposal is not a no-op on the tail: the
fold order for a unit's fixings is the recorded observation index~$k$ --- the fixing's
valid-time coordinate, distinct from its transaction-time position, the point at which its
print was appended to the log (Chapter~\ref{ch:picture}, Section~\ref{sec:bitemporal}) ---
not the order in which proposals arrive, so the proposal is inserted at its observation
position~$k$ and the order-dependent accumulator tail $A_{k},A_{k+1},\dots$ is re-sequenced
and recomputed from $k$ forward. \emph{Replay
determinism:} re-running the map on the recorded observations \emph{in observation-index
order} reproduces every recorded statistic bit-for-bit, because barrier and fixing
observation are deterministic --- the data crosses the boundary once, as a recorded
observation, and is never re-read live (Chapter~\ref{ch:marketdata}); a late insertion
recomputes the tail, and the recomputation is bit-exact.

\medskip\noindent\textbf{Episode --- the future (T1): a duplicate is inert; a late firing is
identical.}
FUT-IDX settles at recorded IDX closes $1{,}000$, $990$, $1{,}005$ (multiplier $10$);
cumulative variation margin is $+50-100+150 = 100$ (Chapters~\ref{ch:contracts}
and~\ref{ch:valuation}). \emph{The gate.} The day-2 settlement event is delivered twice, and
delivered late, in a scheduler that reorders it against \emph{unrelated} firings. Reordering
against unrelated firings leaves FUT-IDX's own fixing sequence --- its observation-index
order --- untouched, so the contract, a function of the record, emits the identical
\USD{100} move; the duplicate carries the same cause-derived identifier and commits once at
the door. Were day~2 instead reordered \emph{within} FUT-IDX's own sequence --- arriving
after day~3 --- it would be inserted at its observation position and the variation-margin
tail recomputed from there, not applied as a late no-op. \emph{Design point.}
\emph{Timeliness depends on the untrusted machines; integrity does not --- a duplicate is
inert, a firing reordered against unrelated events is identical, and one reordered within
its unit's own sequence is re-sequenced and its tail recomputed; all are properties the
scheduler tries and fails to break.}

\medskip\noindent\textbf{Episode --- the variance swap (T5): regenerate firing~127.}
VS-1 accrues a sum-of-squared-returns statistic in UnitStatus, an integer at scale $10^{-8}$,
written by each firing for the next (Chapter~\ref{ch:homes}). At the mid-life checkpoint the
accrued statistic is $A_{126} = 1{,}805{,}000$, and the previous recorded close is fixing
$126 = 1{,}005$. Firing~127 reads $A_{126}$ from the record, forms the return against the
recorded closes, and writes $A_{127}$.

\begin{lstlisting}
accrue :: Close -> Close -> Accrued -> Accrued        -- named, versioned, pure
prop_replay127 :: Property
prop_replay127 = accrue close127 1005 1805000 === recorded A_127
-- the recorded closes are deterministic; A_126 is a recorded integer; so any party
-- re-running the same map on the same recorded fixings obtains the identical A_127.

-- replay folds the fixings in recorded OBSERVATION-INDEX order (the sequence k), never
-- arrival order; a late or out-of-order print is inserted at position k and the tail
-- A_k, A_{k+1}, ... is recomputed from k forward:
prop_lateInsertRecomputesTail :: Fixing -> [Fixing] -> Property
prop_lateInsertRecomputesTail late fs =
  foldAccrue (insertAtObsIndex late fs) === foldAccrue (sortByObsIndex (late : fs))
-- the fold is over an ordered integer sequence, so the recomputed tail is bit-exact.
\end{lstlisting}

Replay is a fold in observation-index order, not arrival order: the immutable log's total
order for VS-1's fixings is the recorded fixing sequence, so an out-of-order or late print
is inserted at its position~$k$ and the accrual tail $A_{k},A_{k+1},\dots$ recomputed from
$k$ forward, bit-for-bit, while the downstream daily CCP-cash figures at positions after~$k$
--- the future's variation margin on the same recorded closes --- are corrected by
authorised compensating transactions (the forward-repair discipline above), not left as posted. Firing~127 is not trusted to have been right
because it ran; it is verified by recomputation, and the recomputation is exact because its
every input --- the accrued integer, the recorded closes --- is on the record.
\emph{Economic verification is re-running the map on recorded inputs and comparing;
determinism is what makes the comparison bit-exact.}

\paragraph{Bitemporality: as-known-at and as-of, made executable.}
Determinism above folds a unit's fixings in valid-time order; the bitemporal regime pins the
two axes of Chapter~\ref{ch:picture} (Section~\ref{sec:bitemporal}) and
Theorem~\ref{thm:timetravel} to generated histories. One property per clause of that theorem.

\begin{lstlisting}
-- BITEMP-1 (as-known-at monotonicity, Theorem 14.4(a)): a correction admitted after a
-- prefix never rewrites the view known at that prefix.
prop_bitemp1 :: History -> Prefix -> CorrectiveObs -> Property
prop_bitemp1 hist p c = isBackDated c ==>
  asKnownAt p hist === asKnownAt p (appendAfter p c hist)

-- BITEMP-2 (as-of determinism, Theorem 14.4(b)): the as-of view over the prefix that
-- carries the correction recomputes bit-for-bit from a fresh valid-time fold.
prop_bitemp2 :: History -> Property
prop_bitemp2 hist = hasBackDated hist ==>
  asOf p' h' (replayFromScratch p') === asOf p' h' hist
  where (p', h') = horizonIncluding (lastBackDated hist) hist
\end{lstlisting}

Both share one non-vacuous precondition: the history must contain a \emph{back-dated}
observation --- one recorded at a transaction-time position after its valid-time index~$k$,
with $k$ strictly inside the existing sequence ($k < \text{end}$). A generator that only ever
appended fixings at the tail would satisfy both properties without ever witnessing a
correction. The generator is
therefore built to reach the branch and to prove it did:

\begin{lstlisting}
genBackDated :: [Fixing] -> Gen History
-- draw a late-flag with rate r; when it fires, append one observation at valid-time index
--   k <- choose (1, maxIndex - 1)   -- strictly inside the sequence: a genuine back-dation
-- otherwise append at the tail (the timely path).
prop_bitemp_fires =                      -- the firing witness, asserted, not hoped
  checkCoverage $ cover 1.0 isBackDated "corrective insertion (k < end)" (genBackDated fs)
\end{lstlisting}

When the late-flag fires the generator constructs $k<\text{end}$ by construction, and its
firing witness asserts the branch (\S 2). On every history that clears the precondition,
BITEMP-1 fires and finds the as-known-at view unchanged by the later correction, and BITEMP-2
fires and finds the as-of view recomputed bit-for-bit --- the two axes exercised, not merely
asserted. C-12.1 (v1.2) names both answers and delegates their indexing here; BITEMP-1 and BITEMP-2
exercise that indexing, and entry PARK-2 is closed (Chapter~\ref{ch:scope}).

\paragraph{(iii) The temporal residual: a safety and liveness model.}
The move algebra owns what is true \emph{within} a transaction --- conservation, coverage,
atomicity, all checked at the door. It does not own the ordering \emph{across} transactions:
which interleavings of independent firings are reachable, and whether an open obligation is
eventually discharged, are temporal questions the algebra cannot answer. The one-touch knock
while pledged is where the residual bites, and it carries an explicit model.

\medskip\noindent\textbf{Episode --- the one-touch (T2): the knock--revalue--call--discharge
model.}
OT-1 (Cast, \S\ref{sec:cast}) is pledged to B under agreement~G as in
\S\ref{sec:knockpledged} (collateral value \USD{200{,}000} against exposure
\USD{150{,}000}). OMEGA's official close prints its barrier close, $79.00 \le 80.00$. Four steps
then interleave: the Monitor emits the knock; OT-1's contract triggers the unit and mints the
conditional payment obligation; the pricing layer, reading the triggered state, marks OT-1 at
$0$, so G's valuation watch opens a margin call for \USD{150{,}000}; the payment obligation
discharges \USD{1{,}000{,}000} (\USD{850{,}000} free, \USD{150{,}000} posted to~B), and the
pledge closes only after the marker returns --- an ordering the coverage invariant forces.
These steps are separate admitted transactions, and their interleaving is the temporal
residual. The model is a state machine over the reachable states, with two invariants held
always and one obligation held eventually. What follows is TLA\textsuperscript{+}-style
specification written as \emph{illustration}, not a machine-checked proof: no model checker is
run inside this specification, and the reference implementation together with its
model-checking lies outside it (\S 5 of the project brief). The operators are read plainly ---
$\square$ is ``at every reachable committed state,'' $\rightsquigarrow$ is ``eventually leads
to,'' $\textrm{WF}$ is ``weak fairness on an action,'' and $\forall$ is ``for all.'' Safety and
Liveness below are therefore \emph{asserted} properties, each with an executable check named in
the regime, not theorems certified here:

\begin{lstlisting}
-- TLA+-style ILLUSTRATION (not machine-checked). Operators: [] always, ~> leads-to, WF_ weak fairness.
Coverage(s)      == \A (w,u): postedMass(w,u) =< Max(owned(w,u), 0)  -- mass; the door's invariant
NoDoubleCount(s) == \A u : extinguished(u,s) => price(u,s) = 0       -- an extinguished exposure is not marked
Safety   == [] ( Coverage /\ NoDoubleCount )                        -- SAFETY: asserted at every committed state
Liveness == WF_vars(dischargeObligation)                            -- fairness hypothesis, INSIDE the formula
              => ( (call = "open") ~> (call = "discharged" \/ call = "closed_out") )  -- LIVENESS: open ~> resolved
-- The condition under which Liveness holds is the WF_ conjunct above, not a sentence beside
-- the formula: drop weak-fairness on dischargeObligation and the leads-to need not hold.
-- Sufficiency (collateral value >= exposure) is NOT a [] invariant: it is lawfully false
-- on the interval [trigger, discharge); only Liveness constrains it. Asserting it as an
-- invariant would be a false theorem, and would forbid a lawful state.
\end{lstlisting}

\emph{Safety} says no reachable committed state double-counts value or violates coverage.
The clause quantifies over units: for every unit whose exposure an admitted lifecycle event
has extinguished, the mark is $0$. A triggered OT-1 is one instance --- no committed state
shows both the option's positive mark and the \USD{1{,}000{,}000} cash it resolves into. Coverage
is excluded because it is checked at the door on every one of the interleaved transactions,
and the marker-before-extinguishment ordering of the final step is coverage refusing to
strip owned while posted mass remains. Safety is not asserted in prose alone: \texttt{NoDoubleCount} is exercised by
\texttt{prop\_noDoubleCount} below, witnessed firing on both extinguishing kinds, and
\texttt{Coverage} is the door invariant every one of the interleaved transactions must clear.
\emph{Liveness} says the margin call opened at the revalue is eventually discharged or closed
out --- the obligation-liveness requirement of the Event Monitor (\S 14.1,
Chapter~\ref{ch:requirements}), asserted as a leads-to property and exercised over generated
histories as a bounded-horizon discharge check: every opened call reaches a resolved state
before its history ends. The fairness under which it holds is the weak-fairness conjunct
written \emph{into} the Liveness formula above --- the Events Executor eventually delivers each
armed firing --- not a hope stated beside it; drop that conjunct and the leads-to need not
hold. The central discipline is what the model \emph{refuses} to assert: between trigger and
discharge the book is lawfully under-collateralised, so sufficiency appears only under
$\rightsquigarrow$, never under $\square$. Coverage --- the mass guard --- and sufficiency ---
the value guard --- are the two guards of Chapter~\ref{ch:collateral}, and here the temporal
model is exactly where their difference is drawn: one is a safety invariant asserted always,
the other a liveness obligation asserted eventually, and a model that confused them would be
unsound in the direction that either fabricates breaches or hides them.

\paragraph{NoDoubleCount fires over every extinguishing event.}
\texttt{NoDoubleCount} quantifies over units, so it must be witnessed on every extinguishing
kind. Two kinds extinguish: OT-1's barrier touch (a triggered one-touch prices at $0$) and
the settled-to-market future's daily settlement (the exposure is owned cash; marking against
the futures level would double-count, \S\ref{sec:pricing-state}). The
collateralised-to-market future is the honest contrast: its return obligation survives, it
is \emph{not} extinguished, and it prices at the last settled level --- the generator must
produce it and the property must not zero it. The
signatures below are Haskell illustration --- the executable form of \texttt{NoDoubleCount} and
its firing regime, kept separate from the TLA\textsuperscript{+}-style model above so the two
notations do not run together in one block:

\begin{lstlisting}
extinguished :: Unit -> State -> Bool                 -- an admitted lifecycle event has extinguished u
extinguished u s =  triggered u s                     -- OT-1: barrier touched, claim resolved
                 || (isFuture u && stm u && vmApplied u s)  -- STM future: day's VM settled, exposure gone
                 -- a CTM future is never extinguished: its return-obligation unit stays live

prop_noDoubleCount :: History -> Property             -- Safety 14.x, quantified over valued units
prop_noDoubleCount h = conjoin
  [ extinguished u (fold h) ==> price u (fold h) === 0 | u <- valuedUnits h ]

prop_noDoubleCount_fires = forAll genMixedHistory $ \h ->  -- witnesses each kind, asserted not hoped
  checkCoverage
  $ cover 1.0 (hasTriggeredOneTouch h) "OT-1 triggered"
  $ cover 1.0 (hasSettledSTMFuture  h) "STM future, VM applied"
  $ cover 1.0 (hasMarkedCTMFuture   h) "CTM future, obligation survives"
  $ prop_noDoubleCount h
\end{lstlisting}

\texttt{genMixedHistory} interleaves three constructions: a knock that triggers a one-touch, a
full variation-margin day applied to a future whose declared regime is settled-to-market, and
the same day applied to a future whose declared regime is collateralised-to-market. Each of the
three labels carries its own firing witness (\S 2), so
\texttt{prop\_noDoubleCount} is exercised on \emph{both} extinguishing kinds --- the OT-1 trigger
and the settled-to-market future's settlement --- and the collateralised-to-market future is
exercised as a non-extinguished unit that must price at its last settled level, never at $0$.
Because \texttt{extinguished} is an
extensible predicate rather than a property that holds by construction, each new extinguishing
lifecycle event-kind must extend \emph{both} the pricing rule (Definition~7.1,
\S\ref{sec:pricing-state}) \emph{and} the \texttt{extinguished} predicate together, or
\texttt{prop\_noDoubleCount} silently under-quantifies and stops witnessing the new kind.

\paragraph{Available free mass is non-negative for every wallet, shorts included.}
The report's available-mass figure (Chapter~\ref{ch:reporting}) carries the same
$\max(\cdot,0)$ as the coverage invariant; without it the figure goes negative for a
negative-owned wallet --- a plain short, or an issuer wallet conservation creates. The
property asserts non-negativity on exactly those wallets.

\begin{lstlisting}
prop_availabilityNonNegative :: History -> Wallet -> Unit -> Property
prop_availabilityNonNegative h w u =
  admitted h ==>
       availableMass h w u >= 0                                  -- the CH.12 REPORT figure, non-negative
  .&&. availableMass h w u
         === max (ownedBal h w u) 0 - sum [ postedG h w u g | g <- agreements h ]
  -- availableMass is the ch12 availability PROJECTION under test, read from the report, NOT
  -- re-derived here: a report shipping (owned - Sigma posted) WITHOUT the max returns < 0 for a
  -- negative-owned wallet and fires this red. Distinct from prop_coverageSignConvention, which
  -- tests the DOOR invariant and would stay green under exactly that broken report.

genNegativeOwned :: Gen History  -- draw a short-flag with rate r; when it fires, build a wallet
  -- with owned < 0 (a plain short / issuer wallet) and read its available-mass report figure.
prop_availability_fires =        -- the firing witness, asserted, not hoped
  checkCoverage $ cover 1.0 hasNegativeOwnedWallet
    "a negative-owned (short/issuer) wallet in the availability report" genNegativeOwned
\end{lstlisting}

The negative-owned branch is a genuine precondition: a generator that only ever built long wallets
would pass \texttt{prop\_availabilityNonNegative} without ever testing the $\max$; its firing
witness asserts the branch (\S 2). The property reads the Chapter~\ref{ch:reporting} availability
figure itself, not a re-derivation of it: on every negative-owned wallet that figure is zero, never
negative, and a report that dropped the $\max$ --- shipping $\mathrm{bal}_{\text{owned}}-\sum_G
\mathrm{bal}_{\text{posted},G}$ --- returns a negative number and fails the property, exactly where
\texttt{prop\_coverageSignConvention} (the door invariant) would still pass. That separation is the
point: the door bound and the report figure are two statements, and this test pins the report.

\paragraph{The fail transition has exactly one writer, and no orphan.}
The walk \texttt{instructed}$\to$\texttt{failed} is written by the settlement-obligation unit's
own contract, on one of the two recorded triggers of Chapter~\ref{ch:settlement}
(Invariant~\ref{inv:one-writer}). The property asserts the writer and forbids an orphan
failed node.

\begin{lstlisting}
prop_failWriterAttribution :: History -> Property
prop_failWriterAttribution h = reachesFailedNode h ==> conjoin
  [      writerOf t === settlementObligationContract h       -- one writer, named
    .&&. triggerOf t `elem` [FailNotice, DueDateNoConfirm]   -- one of two recorded triggers
  | t <- failTransitions h ]  -- EVERY failed node, so a second unit's orphan cannot escape
  -- no `failed' node exists without one of the two triggers on the record.

-- genSettlementHistory (above) already walks instructed -> failed under a fail-flag; the writer
-- attribution and trigger presence are asserted on that same branch.
prop_failWriter_fires =          -- the firing witness, asserted, not hoped
  checkCoverage $ cover 1.0 reachesFailedNode
    "instructed -> failed written by the unit's contract on a recorded trigger" genSettlementHistory
\end{lstlisting}

On every history that reaches the failed node, the transition is authored by the
settlement-obligation unit's contract and carries one of the two triggers on the record; a failed
node written by any other machine, or with no trigger, fails the property --- the writerless fail the
one-writer rule forbids.

\paragraph{Venue is a custody wallet, not a coordinate: the mark-to-market fold.}
Sub-custody is a wallet partition (Chapter~\ref{ch:valuation}); the mark-to-market NAV is a
fold over custody wallets, each priced at its venue's official close. The property asserts
that fold and the absence of any venue axis, and it bites when two or more custody wallets
hold the same unit at different closes.

\begin{lstlisting}
prop_navMkFoldsOverCustodyWallets :: History -> Unit -> Property
prop_navMkFoldsOverCustodyWallets h u = twoCustodyVenues h u ==>
       navMk h u === sum [ closeAt (venueOf w) * ownedBal h w u | w <- custodyWallets h u ]
  .&&. balanceCoordinates h === [Owned, Lent, Posted]  -- NO venue-indexed coordinate exists
  -- venue is WHICH custody wallet, an existing dimension; the position cannot diverge because
  -- there is no venue axis to hold a divergent figure. balanceCoordinates reads the balance
  -- SCHEMA, independent of the fold, so this conjunct is not x===x: add a venue coordinate and
  -- it fires red.

genMultiVenue :: Gen History  -- draw a multi-venue flag with rate r; when it fires, hold u across
  -- >=2 custody wallets (W_A,W_B,...) whose venues publish DIFFERENT official closes.
prop_navMk_fires =            -- the firing witness, asserted, not hoped
  checkCoverage $ cover 1.0 twoCustodyVenues
    "one unit held across >=2 custody wallets at different venue closes" genMultiVenue
\end{lstlisting}

When the flag fires the holding is split across custody wallets at differing closes by construction,
so the cross-venue basis is non-zero and the fold is exercised; its firing witness asserts the
branch (\S 2). The property fails on the
design that adds a venue coordinate: the position could then diverge from the wallet-partition sum,
and \texttt{prop\_navMkFoldsOverCustodyWallets} reports the divergence the single position record
forbids.

\paragraph{Variance fixings read one-plus-closes, and the off-by-one it refuses.}
A realised-variance fold over $N$ daily log-returns reads $N+1$ official closes ---
$C_0$ at the effective date, then $C_1,\dots,C_{N}$ (Chapter~\ref{ch:cdm}). The property
asserts the count relation, that $C_0$ is on the record before fixing 1, and that the fold
sums exactly $N$ squared returns; every generated variance swap witnesses it.

\begin{lstlisting}
prop_varianceFixingsReadOnePlusCloses :: History -> Property
prop_varianceFixingsReadOnePlusCloses h = isVarianceSwap h ==>
       length (closesRead h)  === nFixings h + 1     -- N returns read N+1 closes
  .&&. referenceClose h `recordedBefore` fixing 1 h  -- C_0 on the record before fixing 1
  .&&. length (returns h)     === nFixings h          -- exactly N returns, none missing a close
  .&&. length (foldSquares h) === nFixings h          -- fold sums exactly N squared returns
  -- nFixings h is the DECLARED schedule count from the product terms (the scheduled 252), read
  -- independently of the fold -- NOT length (returns h) -- so these are genuine count checks,
  -- not x===x: an implementation reading N closes for N fixings fails the first conjunct.

genVarianceSwap :: Gen History
-- draw N in [1..252] so N VARIES; fix a reference close C_0 at the effective date,
-- then N daily closes C_1..C_N; fixing i is the return ln(C_i / C_{i-1}).

prop_varFixings_fires =                     -- the firing witness, asserted, not hoped
  checkCoverage $ cover 1.0 isVarianceSwap
    "a variance-swap history: N returns folded over N+1 closes" genVarianceSwap
\end{lstlisting}

Because $C_0$ is drawn at the effective date and $C_1,\dots,C_N$ one per fixing
date, every generated variance-swap history clears the precondition by construction,
and its firing witness asserts the branch (\S 2). The property fails exactly on the off-by-one
design: read $N$ closes
instead of $N+1$ and $C_0$ goes unnamed, the fold produces $N-1$ returns, and
\texttt{prop\_varianceFixingsReadOnePlusCloses} reports the missing close the
reference exists to supply. For $N = 252$ the fold reads $253$ closes and sums $252$
squared returns.

\paragraph{The record is the log: replay reconstructs every observation.}
Every observation enters the log as a moveless observation-recording transaction
(Chapter~\ref{ch:marketdata}), so re-folding a log prefix reconstructs every recorded
observation --- there is no second store. The property discards everything but the log,
rebuilds the whole record, and compares; a generator recording observations outside the log
--- a second door --- fails it at once.

\begin{lstlisting}
prop_replayReconstructsRecord :: History -> Property
prop_replayReconstructsRecord h = hasObservations h ==>
       observations (replayFromLog (logOf h)) === observations h  -- every obs rebuilt from the log
  .&&. record       (replayFromLog (logOf h)) === record h        -- and the whole record with it
  -- observations h / record h are the generator's INDEPENDENT oracle -- what it intended to
  -- record -- NOT defined as replayFromLog; recording an observation outside the log makes
  -- replayFromLog drop it and fires this red (oracle-vs-replay, not x===x). This is the decisive
  -- check for D4: it is the record/log collapse made falsifiable.

genObservationHistory :: Gen History  -- draw histories that record >=1 observation (a close, a
  -- fixing, a barrier print) as a moveless observation-recording transaction on the one log.
prop_replayRecord_fires =              -- the firing witness, asserted, not hoped
  checkCoverage $ cover 1.0 hasObservations
    "a history recording >=1 observation as a moveless transaction" genObservationHistory
\end{lstlisting}

Every generated history records at least one observation through the one door, so the precondition
clears by construction and its firing witness asserts the branch (\S 2). The property fails exactly
on the second-door design: record an observation
anywhere but the log and \texttt{replayFromLog} cannot rebuild it, and
\texttt{prop\_replayReconstructsRecord} reports the store the one-door discipline exists to abolish.

\paragraph{The event door: one registry, mandatory routing, and the quarantine that never drops.}
Every arriving external event is recorded, or its refusal recorded; a kind registers only with
routing defined for it, and an event of an unregistered kind is parked as a \emph{quarantined
event} --- never dropped, never a bare value (C-4.12). Four properties pin the door.

\begin{lstlisting}
prop_registryTotality :: Registration -> Registry -> Property
prop_registryTotality r reg =
      lacksRouter r ==> refusedAtDoor r                                -- router-less: refused
  .&&. conjoin [ isJust (routerFor reg k) | k <- registeredKinds reg ] -- admitted: routing total
  -- the refusal is the DOOR's, not the constructor's (not x===x).
genRegistration :: Gen Registration  -- draw a router-flag rate r: when it fires OMIT the router.
prop_registryTotality_fires =
  checkCoverage $ cover 1.0 lacksRouter "registration without a router -> refused" genRegistration
prop_routingDeterminism :: History -> Property
prop_routingDeterminism h = hasCapturedEvents h ==>
  routeAll (replayFromLog (logOf h)) === recordedRouting h
  -- replay-vs-recorded routing: an independent oracle (not x===x).
genCapturedHistory :: Gen History  -- capture >=1 event of a REGISTERED kind; replay the prefix twice.
prop_routingDeterminism_fires =
  checkCoverage $ cover 1.0 hasCapturedEvents "a captured event routed on replay" genCapturedHistory
prop_quarantineNeverDrop :: History -> Property
prop_quarantineNeverDrop h =
      length (arrivals h) === length (recordedOrQuarantined h)         -- nothing dropped
  .&&. all (\q -> hasProvenance q && hasOpenItem q) (quarantined h)    -- each visible, deadline-bearing
  -- arrivals h is the generator's INDEPENDENT delivered-count, never read back from the log.
genUnregisteredArrivals :: Gen History  -- unregistered-kind flag rate r: deliver an off-registry kind.
prop_quarantineNeverDrop_fires =
  checkCoverage $ cover 1.0 hasUnregisteredArrival
    "unregistered-kind event delivered -> recorded as a quarantined event" genUnregisteredArrivals
prop_lateRegistrationReplay :: History -> Property
prop_lateRegistrationReplay h = registersAfterQuarantine h ==>
      fullState (replayFromLog (logOf h)) === fullState (fold h)       -- replay bit-identical
  .&&. asKnownAt (beforeRegistration h) h === showsQuarantined h       -- pre-reg view still parks it
  -- replay-vs-fold (not x===x); parked processing keyed on the quarantined event's cause.
genLateRegistration :: Gen History  -- late-reg flag rate r: quarantine an off-registry event,
  -- THEN register its kind+router, THEN process the parked event (cause = the quarantined
  -- event), delivered TWICE (idempotence).
prop_lateRegistrationReplay_fires =
  checkCoverage $ cover 1.0 registersAfterQuarantine
    "kind registered after quarantine -> parked event processed once, replay bit-identical" genLateRegistration
\end{lstlisting}

Each firing witness asserts its branch (\S 2): drop the routing-total conjunct and a
router-less kind admits; drop the arrival count and a silent drop passes; key the parked
processing on arrival rather than the quarantined event's cause and replay diverges.

\subsection{What an auditor checks}

A candidate implementation satisfies this chapter when the checks above run green over the
generator universe, each firing witness clearing the $\ge 1\%$ floor (\S 2): the structural
invariants of Chapter~\ref{ch:invariants} hold as
invariants and no adversarial history reaches a forbidden state; the economic oracle
reproduces every recorded transaction from its recorded event, and a seeded divergence is
repaired only by an authorised compensating transaction, never by an edit; the four product-graph properties confirm that
watches, state schema, and actions never disagree; across a generated split of a referenced
underlying the derivative's payoff is invariant and no stale-barrier knock fires; moves are the
sole balance mutator; a duplicate firing is inert and a late firing folds in at its observation
index --- identical when reordered against unrelated events, its tail recomputed when reordered
within its unit's own sequence; firing~127 recomputes bit-for-bit; conservation holds per (unit,
coordinate, agreement) and the cross-agreement netting branch --- post under $G_1$, receive
under $G_2$, aggregate posted axis zero --- reads the wallet encumbered where a net read would
show zero; two dividends in flight are both live in the position-state collection and each
retires on its own payment date; a coordinated cascade delivered twice commits each of its $n$
legs exactly once and all $n$; the bitemporal pair BITEMP-1 and BITEMP-2 hold with the
back-dated-observation branch ($k<\text{end}$) exercised; every arriving external event is
recorded or refused, no event kind registers without routing defined for it, an event of an
unregistered kind is parked as a quarantined event with provenance and an open item rather than
dropped, and a kind registered after quarantine processes its parked event once with replay
bit-identical; and the
knock--revalue--call--discharge model holds its two safety invariants
always and its liveness obligation eventually, with sufficiency asserted only as the latter.
Each is mechanical, each is replayable from a seed, and each fails closed. A property that
cannot be checked this way is redesigned until it can be: property-based testing is a
construction principle here, not a phase at the end.
